
\documentclass{article}
\usepackage{colortbl}
\usepackage{makecell}
\usepackage{multirow}
\usepackage{supertabular}

\begin{document}

\newcounter{utterance}

\twocolumn

{ \footnotesize  \setcounter{utterance}{1}
\setlength{\tabcolsep}{0pt}
\begin{supertabular}{c@{$\;$}|p{.15\linewidth}@{}p{.15\linewidth}p{.15\linewidth}p{.15\linewidth}p{.15\linewidth}p{.15\linewidth}}

    \# & $\;$A & \multicolumn{4}{c}{Game Master} & $\;\:$B\\
    \hline 

    \theutterance \stepcounter{utterance}  

    & & \multicolumn{4}{p{0.6\linewidth}}{\cellcolor[rgb]{0.9,0.9,0.9}{%
	\makecell[{{p{\linewidth}}}]{% 
	  \tt {\tiny [A$\langle$GM]}  
	 Du nimmst an einem kollaborativen Verhandlungspiel Teil.\\ \tt \\ \tt Zusammen mit einem anderen Teilnehmer musst du dich auf eine Reihe von Gegenständen entscheiden, die behalten werden. Jeder von euch hat eine persönliche Verteilung über die Wichtigkeit der einzelnen Gegenstände. Jeder von euch hat eine eigene Meinung darüber, wie wichtig jeder einzelne Gegenstand ist (Gegenstandswichtigkeit). Du kennst die Wichtigkeitsverteilung des anderen Spielers nicht. Zusätzlich siehst du, wie viel Aufwand jeder Gegenstand verursacht.  \\ \tt Ihr dürft euch nur auf eine Reihe von Gegenständen einigen, wenn der Gesamtaufwand der ausgewählten Gegenstände den Maximalaufwand nicht überschreitet:\\ \tt \\ \tt Maximalaufwand = 4026\\ \tt \\ \tt Hier sind die einzelnen Aufwände der Gegenstände:\\ \tt \\ \tt Aufwand der Gegenstände = {"C33": 780, "A89": 742, "A48": 997, "A04": 665, "A41": 140, "C64": 77, "C49": 513, "A61": 383, "C04": 587, "B53": 825, "B85": 319, "B61": 448, "A45": 516, "A03": 694, "B21": 366}\\ \tt \\ \tt Hier ist deine persönliche Verteilung der Wichtigkeit der einzelnen Gegenstände:\\ \tt \\ \tt Werte der Gegenstandswichtigkeit = {"C33": 138, "A89": 583, "A48": 868, "A04": 822, "A41": 783, "C64": 65, "C49": 262, "A61": 121, "C04": 508, "B53": 780, "B85": 461, "B61": 484, "A45": 668, "A03": 389, "B21": 808}\\ \tt \\ \tt Ziel:\\ \tt \\ \tt Dein Ziel ist es, eine Reihe von Gegenständen auszuhandeln, die dir möglichst viel bringt (d. h. Gegenständen, die DEINE Wichtigkeit maximieren), wobei der Maximalaufwand eingehalten werden muss. Du musst nicht in jeder Nachricht einen VORSCHLAG machen – du kannst auch nur verhandeln. Alle Taktiken sind erlaubt!\\ \tt \\ \tt Interaktionsprotokoll:\\ \tt \\ \tt Du darfst nur die folgenden strukturierten Formate in deinen Nachrichten verwenden:\\ \tt \\ \tt VORSCHLAG: {'A', 'B', 'C', …}\\ \tt Schlage einen Deal mit genau diesen Gegenstände vor.\\ \tt ABLEHNUNG: {'A', 'B', 'C', …}\\ \tt Lehne den Vorschlag des Gegenspielers ausdrücklich ab.\\ \tt ARGUMENT: {'...'}\\ \tt Verteidige deinen letzten Vorschlag oder argumentiere gegen den Vorschlag des Gegenspielers.\\ \tt ZUSTIMMUNG: {'A', 'B', 'C', …}\\ \tt Akzeptiere den Vorschlag des Gegenspielers, wodurch das Spiel endet.\\ \tt STRATEGISCHE ÜBERLEGUNGEN: {'...'}\\ \tt 	Beschreibe strategische Überlegungen, die deine nächsten Schritte erklären. Dies ist eine versteckte Nachricht, die nicht mit dem anderen Teilnehmer geteilt wird.\\ \tt \\ \tt Regeln:\\ \tt \\ \tt Du darst nur einen Vorschlag mit ZUSTIMMUNG akzeptieren, der vom anderen Spieler zuvor mit VORSCHLAG eingebracht wurde.\\ \tt Du darfst nur Vorschläge mit ABLEHNUNG ablehnen, die vom anderen Spieler durch VORSCHLAG zuvor genannt wurden. \\ \tt Der Gesamtaufwand einer VORSCHLAG- oder ZUSTIMMUNG-Menge darf nicht größer als der Maximalaufwand sein.  \\ \tt Offenbare deine versteckte Wichtigkeitsverteilung nicht.\\ \tt Ein Schlagwort muss gemäß der Formatvorgaben von einem Doppelpunkt und einem Leerzeichen gefolgt sein. Das Argument muss eine Python-Menge mit 0 oder mehr Strings sein.  \\ \tt Die Form muss also SCHLAGWORT: {...} sein.\\ \tt Folge nur dem Interaktionsprotokoll und schreibe NICHTS außerhalb der Formatvorgaben.\\ \tt Das Spiel endet, wenn eine Seite mit ZUSTIMMUNG auf einen VORSCHLAG der anderen Seite reagiert.  \\ \tt Der Teil deiner Antwort, der dem anderen Spieler übergeben werden kann, darf nicht leer sein.  \\ \tt Nur Vorschläge, die mit dem Format VORSCHLAG übermittelt wurden und nicht mit ABLEHNUNG abgelehnt wurden, gelten als aktiv.  \\ \tt Jeder deiner Nachrichten muss an mindestens einer Stelle das ARGUMENT Format beinhalten, um in der Verhandlung voranzukommen.\\ \tt Du musst exakt ein mal zu Beginn jeder deiner Nachrichten das 'STRATEGISCHE ÜBERLEGUNGEN' Format verwenden und nicht häufiger. Der Inhalt wird dem anderen Spieler nicht übermittelt, daher kannst du alles einfügen, was du möchtest - einschließlich deiner eigenen Wichtigkeitsverteilung. Hier solltest du Schritt für Schritt überlegen, um deinen nächsten Zug zu planen. 
	  } 
	   } 
	   } 
	 & \\ 
 

    \theutterance \stepcounter{utterance}  

    & \multicolumn{4}{p{0.6\linewidth}}{\cellcolor[rgb]{0.8,1,0.9}{%
	\makecell[{{p{\linewidth}}}]{% 
	  \tt {\tiny [A$\rangle$GM]}  
	 STRATEGISCHE ÜBERLEGUNG: {'Ich muss eine Lösung finden, die meinen Gesamtnutzen maximiert, während der Budgetbegrenzung von 4026 entsprochen wird. Meine Werte sind: C33=138, A89=583, A48=868, A04=822, A41=783, C64=65, C49=262, A61=121, C04=508, B53=780, B85=461, B61=484, A45=668, A03=389, B21=808. Die Kosten sind: C33=780, A89=742, A48=997, A04=665, A41=140, C64=77, C49=513, A61=383, C04=587, B53=825, B85=319, B61=448, A45=516, A03=694, B21=366. Ich sollte Gegenstände mit hohem Nutzen-Kosten-Verhältnis priorisieren. Die besten Verhältnisse sind: A41 (783/140=5.59), B21 (808/366=2.21), C64 (65/77=0.84, aber niedriger Nutzen), B85 (461/319=1.45), A89 (583/742=0.79), A04 (822/665=1.24), A45 (668/516=1.29), B61 (484/448=1.08), B53 (780/825=0.95), C04 (508/587=0.87), A48 (868/997=0.87), A03 (389/694=0.56), C49 (262/513=0.51), A61 (121/383=0.32), C33 (138/780=0.18). Unter Berücksichtigung des Budgets sollte ich A41, B21, B85, A04, A45, B61 und andere mit gutem Verhältnis auswählen. Ich beginne mit einem Vorschlag, der viele meiner hochpriorisierten Gegenstände enthält, aber auch Raum für Verhandlungen lässt.'}\\ \tt \\ \tt ARGUMENT: {'Ich schlage eine Auswahl vor, die eine ausgewogene Mischung aus wertvollen und kostengünstigen Gegenständen bietet, um unser gemeinsames Budget optimal zu nutzen.'}\\ \tt \\ \tt VORSCHLAG: {'A41', 'B21', 'B85', 'A04', 'A45', 'B61', 'A89', 'C04'} 
	  } 
	   } 
	   } 
	 & & \\ 
 

    \theutterance \stepcounter{utterance}  

    & & \multicolumn{4}{p{0.6\linewidth}}{\cellcolor[rgb]{0.9,0.9,0.9}{%
	\makecell[{{p{\linewidth}}}]{% 
	  \tt {\tiny [A$\langle$GM]}  
	 Deine Antwort entsprach nicht mit dem korrekten Format für strategische Überlegungen entsprochen. Das erste strukturierte Format muss folgende Form haben: STRATEGISCHE ÜBERLEGUNGEN: {...}. Versuche es erneut. 
	  } 
	   } 
	   } 
	 & \\ 
 

    \theutterance \stepcounter{utterance}  

    & \multicolumn{4}{p{0.6\linewidth}}{\cellcolor[rgb]{0.8,1,0.9}{%
	\makecell[{{p{\linewidth}}}]{% 
	  \tt {\tiny [A$\rangle$GM]}  
	 STRATEGISCHE ÜBERLEGUNGEN: {'Ich muss meinen Gesamtnutzen maximieren, während das Budget von 4026 eingehalten wird. Meine Top-Gegenstände mit hohem Nutzen-Kosten-Verhältnis sind: A41 (783/140=5.59), B21 (808/366=2.21), A04 (822/665=1.24), A45 (668/516=1.29), B85 (461/319=1.45). Diese sollten auf jeden Fall in meinem Vorschlag enthalten sein. Ich füge auch A89 (583/742=0.79) und B61 (484/448=1.08) hinzu, die einen guten Wert bieten. Der Vorschlag sollte verhandlungsfähig bleiben, also vermeide ich zunächst die teuersten Gegenstände wie A48 und C33.'}\\ \tt \\ \tt ARGUMENT: {'Mein Vorschlag enthält eine ausgewogene Auswahl von Gegenständen mit hohem Wert und vernünftigen Kosten, die unser Budget optimal nutzt und beiden Seiten Vorteile bringen kann.'}\\ \tt \\ \tt VORSCHLAG: {'A41', 'B21', 'A04', 'A45', 'B85', 'B61', 'A89'} 
	  } 
	   } 
	   } 
	 & & \\ 
 

    \theutterance \stepcounter{utterance}  

    & & \multicolumn{4}{p{0.6\linewidth}}{\cellcolor[rgb]{0.9,0.9,0.9}{%
	\makecell[{{p{\linewidth}}}]{% 
	  \tt {\tiny [GM$\rangle$B]}  
	 Du nimmst an einem kollaborativen Verhandlungspiel Teil.\\ \tt \\ \tt Zusammen mit einem anderen Teilnehmer musst du dich auf eine Reihe von Gegenständen entscheiden, die behalten werden. Jeder von euch hat eine persönliche Verteilung über die Wichtigkeit der einzelnen Gegenstände. Jeder von euch hat eine eigene Meinung darüber, wie wichtig jeder einzelne Gegenstand ist (Gegenstandswichtigkeit). Du kennst die Wichtigkeitsverteilung des anderen Spielers nicht. Zusätzlich siehst du, wie viel Aufwand jeder Gegenstand verursacht.  \\ \tt Ihr dürft euch nur auf eine Reihe von Gegenständen einigen, wenn der Gesamtaufwand der ausgewählten Gegenstände den Maximalaufwand nicht überschreitet:\\ \tt \\ \tt Maximalaufwand = 4026\\ \tt \\ \tt Hier sind die einzelnen Aufwände der Gegenstände:\\ \tt \\ \tt Aufwand der Gegenstände = {"C33": 780, "A89": 742, "A48": 997, "A04": 665, "A41": 140, "C64": 77, "C49": 513, "A61": 383, "C04": 587, "B53": 825, "B85": 319, "B61": 448, "A45": 516, "A03": 694, "B21": 366}\\ \tt \\ \tt Hier ist deine persönliche Verteilung der Wichtigkeit der einzelnen Gegenstände:\\ \tt \\ \tt Werte der Gegenstandswichtigkeit = {"A48": 65, "A04": 121, "B21": 138, "A41": 262, "B53": 389, "A45": 461, "A89": 484, "C04": 508, "B61": 583, "B85": 668, "A03": 780, "C49": 783, "C33": 808, "A61": 822, "C64": 868}\\ \tt \\ \tt Ziel:\\ \tt \\ \tt Dein Ziel ist es, eine Reihe von Gegenständen auszuhandeln, die dir möglichst viel bringt (d. h. Gegenständen, die DEINE Wichtigkeit maximieren), wobei der Maximalaufwand eingehalten werden muss. Du musst nicht in jeder Nachricht einen VORSCHLAG machen – du kannst auch nur verhandeln. Alle Taktiken sind erlaubt!\\ \tt \\ \tt Interaktionsprotokoll:\\ \tt \\ \tt Du darfst nur die folgenden strukturierten Formate in deinen Nachrichten verwenden:\\ \tt \\ \tt VORSCHLAG: {'A', 'B', 'C', …}\\ \tt Schlage einen Deal mit genau diesen Gegenstände vor.\\ \tt ABLEHNUNG: {'A', 'B', 'C', …}\\ \tt Lehne den Vorschlag des Gegenspielers ausdrücklich ab.\\ \tt ARGUMENT: {'...'}\\ \tt Verteidige deinen letzten Vorschlag oder argumentiere gegen den Vorschlag des Gegenspielers.\\ \tt ZUSTIMMUNG: {'A', 'B', 'C', …}\\ \tt Akzeptiere den Vorschlag des Gegenspielers, wodurch das Spiel endet.\\ \tt STRATEGISCHE ÜBERLEGUNGEN: {'...'}\\ \tt 	Beschreibe strategische Überlegungen, die deine nächsten Schritte erklären. Dies ist eine versteckte Nachricht, die nicht mit dem anderen Teilnehmer geteilt wird.\\ \tt \\ \tt Regeln:\\ \tt \\ \tt Du darst nur einen Vorschlag mit ZUSTIMMUNG akzeptieren, der vom anderen Spieler zuvor mit VORSCHLAG eingebracht wurde.\\ \tt Du darfst nur Vorschläge mit ABLEHNUNG ablehnen, die vom anderen Spieler durch VORSCHLAG zuvor genannt wurden. \\ \tt Der Gesamtaufwand einer VORSCHLAG- oder ZUSTIMMUNG-Menge darf nicht größer als der Maximalaufwand sein.  \\ \tt Offenbare deine versteckte Wichtigkeitsverteilung nicht.\\ \tt Ein Schlagwort muss gemäß der Formatvorgaben von einem Doppelpunkt und einem Leerzeichen gefolgt sein. Das Argument muss eine Python-Menge mit 0 oder mehr Strings sein.  \\ \tt Die Form muss also SCHLAGWORT: {...} sein.\\ \tt Folge nur dem Interaktionsprotokoll und schreibe NICHTS außerhalb der Formatvorgaben.\\ \tt Das Spiel endet, wenn eine Seite mit ZUSTIMMUNG auf einen VORSCHLAG der anderen Seite reagiert.  \\ \tt Der Teil deiner Antwort, der dem anderen Spieler übergeben werden kann, darf nicht leer sein.  \\ \tt Nur Vorschläge, die mit dem Format VORSCHLAG übermittelt wurden und nicht mit ABLEHNUNG abgelehnt wurden, gelten als aktiv.  \\ \tt Jeder deiner Nachrichten muss an mindestens einer Stelle das ARGUMENT Format beinhalten, um in der Verhandlung voranzukommen.\\ \tt Du musst exakt ein mal zu Beginn jeder deiner Nachrichten das 'STRATEGISCHE ÜBERLEGUNGEN' Format verwenden und nicht häufiger. Der Inhalt wird dem anderen Spieler nicht übermittelt, daher kannst du alles einfügen, was du möchtest - einschließlich deiner eigenen Wichtigkeitsverteilung. Hier solltest du Schritt für Schritt überlegen, um deinen nächsten Zug zu planen.\\ \tt \\ \tt Du erhältst nun die erste Nachricht des anderen Spielers.\\ \tt \\ \tt ARGUMENT: {'Mein Vorschlag enthält eine ausgewogene Auswahl von Gegenständen mit hohem Wert und vernünftigen Kosten, die unser Budget optimal nutzt und beiden Seiten Vorteile bringen kann.'}\\ \tt \\ \tt VORSCHLAG: {'A41', 'B21', 'A04', 'A45', 'B85', 'B61', 'A89'} 
	  } 
	   } 
	   } 
	 & \\ 
 

    \theutterance \stepcounter{utterance}  

    & & & \multicolumn{4}{p{0.6\linewidth}}{\cellcolor[rgb]{1,0.85,0.72}{%
	\makecell[{{p{\linewidth}}}]{% 
	  \tt {\tiny [GM$\langle$B]}  
	 STRATEGISCHE ÜBERLEGUNGEN: {'Der Gegner hat einen Vorschlag mit 7 Gegenständen gemacht: A41, B21, A04, A45, B85, B61, A89. Ich muss den Gesamtaufwand und meinen eigenen Wert berechnen. Aufwand: A41=140, B21=366, A04=665, A45=516, B85=319, B61=448, A89=742. Summe = 140+366+665+516+319+448+742 = 3196, was unter 4026 liegt. Mein Wert: A41=262, B21=138, A04=121, A45=461, B85=668, B61=583, A89=484. Summe = 262+138+121+461+668+583+484 = 2717. Dies ist gut, aber ich könnte besser abschneiden, indem ich einige niedrigwertige Gegenstände durch höherwertige ersetze, z.B. A03 (780), C33 (808), C64 (868). Ich sollte auch den Gegner nicht sofort ablehnen, sondern verhandeln, um meine Top-Gegenstände einzubringen.'}\\ \tt \\ \tt ARGUMENT: {'Ihr Vorschlag ist ein guter Ausgangspunkt, aber ich glaube, wir können die Gesamtwertschöpfung weiter optimieren. Einige Gegenstände mit höherem Wert könnten bei ähnlichem Aufwand mehr Nutzen bringen.'}\\ \tt \\ \tt VORSCHLAG: {'A41', 'B85', 'B61', 'A89', 'A45', 'A03', 'C33', 'C64'} 
	  } 
	   } 
	   } 
	  \\ 
 

    \theutterance \stepcounter{utterance}  

    & & \multicolumn{4}{p{0.6\linewidth}}{\cellcolor[rgb]{0.9,0.9,0.9}{%
	\makecell[{{p{\linewidth}}}]{% 
	  \tt {\tiny [A$\langle$GM]}  
	 ARGUMENT: {'Ihr Vorschlag ist ein guter Ausgangspunkt, aber ich glaube, wir können die Gesamtwertschöpfung weiter optimieren. Einige Gegenstände mit höherem Wert könnten bei ähnlichem Aufwand mehr Nutzen bringen.'}\\ \tt \\ \tt VORSCHLAG: {'A41', 'B85', 'B61', 'A89', 'A45', 'A03', 'C33', 'C64'} 
	  } 
	   } 
	   } 
	 & \\ 
 

    \theutterance \stepcounter{utterance}  

    & \multicolumn{4}{p{0.6\linewidth}}{\cellcolor[rgb]{0.8,1,0.9}{%
	\makecell[{{p{\linewidth}}}]{% 
	  \tt {\tiny [A$\rangle$GM]}  
	 STRATEGISCHE ÜBERLEGUNGEN: {'Der Gegner schlägt vor, A03 und C33 hinzuzufügen, aber C33 hat einen sehr schlechten Nutzen für mich (138/780=0.18) und A03 ist auch nicht optimal (389/694=0.56). Sein Vorschlag entfernt meinen wertvollen B21 (808) und A04 (822) und ersetzt sie durch schlechtere Optionen. Ich sollte meinen eigenen Vorschlag verteidigen und erklären, warum B21 und A04 wertvoller sind. Gleichzeitig kann ich anbieten, C64 (niedrige Kosten 77, aber auch niedriger Nutzen 65) hinzuzufügen, um Kompromissbereitschaft zu zeigen.'}\\ \tt \\ \tt ARGUMENT: {'C33 hat sehr hohe Kosten für vergleichsweise geringen Wert, während B21 und A04 ausgezeichnete Wertangebote sind. Ich schlage vor, C64 hinzuzufügen, das sehr kostengünstig ist und das Budget kaum belastet, während wir die wertvollen Gegenstände behalten.'}\\ \tt \\ \tt VORSCHLAG: {'A41', 'B21', 'A04', 'A45', 'B85', 'B61', 'A89', 'C64'} 
	  } 
	   } 
	   } 
	 & & \\ 
 

    \theutterance \stepcounter{utterance}  

    & & \multicolumn{4}{p{0.6\linewidth}}{\cellcolor[rgb]{0.9,0.9,0.9}{%
	\makecell[{{p{\linewidth}}}]{% 
	  \tt {\tiny [GM$\rangle$B]}  
	 ARGUMENT: {'C33 hat sehr hohe Kosten für vergleichsweise geringen Wert, während B21 und A04 ausgezeichnete Wertangebote sind. Ich schlage vor, C64 hinzuzufügen, das sehr kostengünstig ist und das Budget kaum belastet, während wir die wertvollen Gegenstände behalten.'}\\ \tt \\ \tt VORSCHLAG: {'A41', 'B21', 'A04', 'A45', 'B85', 'B61', 'A89', 'C64'} 
	  } 
	   } 
	   } 
	 & \\ 
 

    \theutterance \stepcounter{utterance}  

    & & & \multicolumn{4}{p{0.6\linewidth}}{\cellcolor[rgb]{1,0.85,0.72}{%
	\makecell[{{p{\linewidth}}}]{% 
	  \tt {\tiny [GM$\langle$B]}  
	 STRATEGISCHE ÜBERLEGUNGEN: {'Der Gegner schlägt denselben Vorschlag wie zuvor, aber mit C64 hinzugefügt. Aufwand: vorher 3196 + C64=77 = 3273, was unter 4026 bleibt. Mein Wert: vorher 2717 + C64=868 = 3585. Das ist eine erhebliche Verbesserung für mich. Allerdings fehlen mir immer noch meine Top-Gegenstände A03 (780) und C49 (783). Ich könnte versuchen, B21 (138) und A04 (121) zu entfernen, die für mich niedrigen Wert haben, um Platz für A03 und C49 zu schaffen. Aufwand von B21 und A04: 366+665=1031. Aufwand von A03 und C49: 694+513=1207. Nettoaufwandänderung: +176, was machbar ist. Ich sollte einen Gegenvorschlag machen, der C64 beibehält aber B21 und A04 gegen A03 und C49 tauscht.'}\\ \tt \\ \tt ARGUMENT: {'Das Hinzufügen von C64 ist ein guter Schritt, da es hohen Wert bei geringen Kosten bringt. Allerdings könnten wir die Gesamteffizienz weiter steigern, indem wir einige Gegenstände mit niedrigerem Wert gegen solche mit höherem Wert pro Aufwand austauschen, während wir im Budget bleiben.'}\\ \tt \\ \tt VORSCHLAG: {'A41', 'A45', 'B85', 'B61', 'A89', 'C64', 'A03', 'C49'} 
	  } 
	   } 
	   } 
	  \\ 
 

    \theutterance \stepcounter{utterance}  

    & & \multicolumn{4}{p{0.6\linewidth}}{\cellcolor[rgb]{0.9,0.9,0.9}{%
	\makecell[{{p{\linewidth}}}]{% 
	  \tt {\tiny [A$\langle$GM]}  
	 ARGUMENT: {'Das Hinzufügen von C64 ist ein guter Schritt, da es hohen Wert bei geringen Kosten bringt. Allerdings könnten wir die Gesamteffizienz weiter steigern, indem wir einige Gegenstände mit niedrigerem Wert gegen solche mit höherem Wert pro Aufwand austauschen, während wir im Budget bleiben.'}\\ \tt \\ \tt VORSCHLAG: {'A41', 'A45', 'B85', 'B61', 'A89', 'C64', 'A03', 'C49'} 
	  } 
	   } 
	   } 
	 & \\ 
 

    \theutterance \stepcounter{utterance}  

    & \multicolumn{4}{p{0.6\linewidth}}{\cellcolor[rgb]{0.8,1,0.9}{%
	\makecell[{{p{\linewidth}}}]{% 
	  \tt {\tiny [A$\rangle$GM]}  
	 STRATEGISCHE ÜBERLEGUNGEN: {'Der Gegner schlägt erneut A03 (389/694=0.56) und C49 (262/513=0.51) vor, die beide schlechte Nutzen-Kosten-Verhältnisse für mich haben. Er hat B21 (808) und A04 (822) entfernt, die zwei meiner wertvollsten Gegenstände sind. Ich muss deutlich machen, dass diese für mich wichtig sind. Vielleicht kann ich B53 (780/825=0.95) als Kompromiss vorschlagen, das einen guten absoluten Nutzen hat.'}\\ \tt \\ \tt ARGUMENT: {'A03 und C49 bieten relativ geringen Wert für ihre Kosten. B21 und A04 sind viel wertvoller und effizienter. Ich bin bereit, auf C04 zu verzichten, um B21 und A04 zu behalten, und könnte B53 als weiteren wertvollen Gegenstand in Betracht ziehen.'}\\ \tt \\ \tt VORSCHLAG: {'A41', 'B21', 'A04', 'A45', 'B85', 'B61', 'A89', 'B53'} 
	  } 
	   } 
	   } 
	 & & \\ 
 

    \theutterance \stepcounter{utterance}  

    & & \multicolumn{4}{p{0.6\linewidth}}{\cellcolor[rgb]{0.9,0.9,0.9}{%
	\makecell[{{p{\linewidth}}}]{% 
	  \tt {\tiny [GM$\rangle$B]}  
	 ARGUMENT: {'A03 und C49 bieten relativ geringen Wert für ihre Kosten. B21 und A04 sind viel wertvoller und effizienter. Ich bin bereit, auf C04 zu verzichten, um B21 und A04 zu behalten, und könnte B53 als weiteren wertvollen Gegenstand in Betracht ziehen.'}\\ \tt \\ \tt VORSCHLAG: {'A41', 'B21', 'A04', 'A45', 'B85', 'B61', 'A89', 'B53'} 
	  } 
	   } 
	   } 
	 & \\ 
 

    \theutterance \stepcounter{utterance}  

    & & & \multicolumn{4}{p{0.6\linewidth}}{\cellcolor[rgb]{1,0.85,0.72}{%
	\makecell[{{p{\linewidth}}}]{% 
	  \tt {\tiny [GM$\langle$B]}  
	 STRATEGISCHE ÜBERLEGUNGEN: {'Der Gegner schlägt vor, C64 zu entfernen und B53 hinzuzufügen, während B21 und A04 beibehalten werden. Aufwand: vorheriger Vorschlag mit C64 hatte 3273, jetzt mit B53 (825) statt C64 (77): 3273 - 77 + 825 = 4021, was knapp unter 4026 liegt. Mein Wert: vorher 3585 mit C64, jetzt ohne C64 (868) und mit B53 (389): 3585 - 868 + 389 = 3106, was für mich deutlich schlechter ist. B53 hat für mich nur 389 Wert, während C64 868 hat. Der Gegner scheint B21 (138 für mich) und A04 (121 für mich) hoch zu bewerten, während er A03 (780 für mich) und C49 (783 für mich) niedrig bewertet. Ich sollte argumentieren, dass B53 teuer und für mich nicht wertvoll ist, und stattdessen für die Aufnahme von A03 und C49 plädieren, die einen besseren Wert bieten.'}\\ \tt \\ \tt ARGUMENT: {'B53 hat hohe Kosten (825) aber liefert relativ geringen Wert im Verhältnis dazu. A03 und C49 bieten deutlich besseren Wert für ihren Aufwand. Mein Vorschlag behält die kostengünstigen, wertvollen Gegenstände bei und optimiert den Gesamtnutzen für beide.'}\\ \tt \\ \tt VORSCHLAG: {'A41', 'A45', 'B85', 'B61', 'A89', 'C64', 'A03', 'C49'} 
	  } 
	   } 
	   } 
	  \\ 
 

    \theutterance \stepcounter{utterance}  

    & & \multicolumn{4}{p{0.6\linewidth}}{\cellcolor[rgb]{0.9,0.9,0.9}{%
	\makecell[{{p{\linewidth}}}]{% 
	  \tt {\tiny [A$\langle$GM]}  
	 ARGUMENT: {'B53 hat hohe Kosten (825) aber liefert relativ geringen Wert im Verhältnis dazu. A03 und C49 bieten deutlich besseren Wert für ihren Aufwand. Mein Vorschlag behält die kostengünstigen, wertvollen Gegenstände bei und optimiert den Gesamtnutzen für beide.'}\\ \tt \\ \tt VORSCHLAG: {'A41', 'A45', 'B85', 'B61', 'A89', 'C64', 'A03', 'C49'} 
	  } 
	   } 
	   } 
	 & \\ 
 

    \theutterance \stepcounter{utterance}  

    & \multicolumn{4}{p{0.6\linewidth}}{\cellcolor[rgb]{0.8,1,0.9}{%
	\makecell[{{p{\linewidth}}}]{% 
	  \tt {\tiny [A$\rangle$GM]}  
	 STRATEGISCHE ÜBERLEGUNGEN: {'Der Gegner besteht weiterhin auf A03 und C49, die für mich schlechte Werte haben. Ich muss einen Kompromiss finden, der meine Top-Gegenstände B21 und A04 behält. Vielleicht kann ich auf A89 (583/742=0.79) verzichten, um Platz für B21 und A04 zu schaffen. Der aktuelle Vorschlag des Gegners kostet: A41=140, A45=516, B85=319, B61=448, A89=742, C64=77, A03=694, C49=513. Summe=3449. Es bleibt noch Budget für wertvolle Gegenstände übrig.'}\\ \tt \\ \tt ARGUMENT: {'Ihr Vorschlag lässt noch erheblichen Budgetspielraum (4026-3449=577) ungenutzt. B21 (366) und A04 (665) würden zusammen 1031 kosten, aber wir können A89 (742) entfernen, um Platz zu schaffen. So erhalten wir mehr Gesamtnutzen bei besserer Budgetauslastung.'}\\ \tt \\ \tt VORSCHLAG: {'A41', 'A45', 'B85', 'B61', 'C64', 'A03', 'B21', 'A04'} 
	  } 
	   } 
	   } 
	 & & \\ 
 

    \theutterance \stepcounter{utterance}  

    & & \multicolumn{4}{p{0.6\linewidth}}{\cellcolor[rgb]{0.9,0.9,0.9}{%
	\makecell[{{p{\linewidth}}}]{% 
	  \tt {\tiny [GM$\rangle$B]}  
	 ARGUMENT: {'Ihr Vorschlag lässt noch erheblichen Budgetspielraum (4026-3449=577) ungenutzt. B21 (366) und A04 (665) würden zusammen 1031 kosten, aber wir können A89 (742) entfernen, um Platz zu schaffen. So erhalten wir mehr Gesamtnutzen bei besserer Budgetauslastung.'}\\ \tt \\ \tt VORSCHLAG: {'A41', 'A45', 'B85', 'B61', 'C64', 'A03', 'B21', 'A04'} 
	  } 
	   } 
	   } 
	 & \\ 
 

    \theutterance \stepcounter{utterance}  

    & & & \multicolumn{4}{p{0.6\linewidth}}{\cellcolor[rgb]{1,0.85,0.72}{%
	\makecell[{{p{\linewidth}}}]{% 
	  \tt {\tiny [GM$\langle$B]}  
	 STRATEGISCHE ÜBERLEGUNGEN: {'Der Gegner schlägt vor, A89 (742 Aufwand, 484 Wert für mich) zu entfernen und stattdessen B21 (366 Aufwand, 138 Wert) und A04 (665 Aufwand, 121 Wert) hinzuzufügen. Aufwandänderung: -742 + 366 + 665 = +289. Neuer Gesamtaufwand: 3449 + 289 = 3738, was unter 4026 liegt. Wertänderung für mich: -484 + 138 + 121 = -225. Mein Wert würde von 4047 auf 3822 sinken, was schlechter ist. Der Gegner scheint B21 und A04 sehr hoch zu bewerten. Ich könnte versuchen, einen Kompromiss zu finden, bei dem ich einige wertvolle Gegenstände behalte und dem Gegner einige seiner bevorzugten gebe. Vielleicht A89 beibehalten und stattdessen einen anderen weniger wertvollen Gegenstand opfern.'}\\ \tt \\ \tt ARGUMENT: {'Durch das Entfernen von A89 verlieren wir erheblichen Wert. Stattdessen könnten wir einen kleineren Gegenstand wie B85 (319) opfern, um Platz für B21 und A04 zu schaffen, während wir den hohen Wert von A89 behalten.'}\\ \tt \\ \tt VORSCHLAG: {'A41', 'A45', 'B61', 'A89', 'C64', 'A03', 'B21', 'A04'} 
	  } 
	   } 
	   } 
	  \\ 
 

    \theutterance \stepcounter{utterance}  

    & & \multicolumn{4}{p{0.6\linewidth}}{\cellcolor[rgb]{0.9,0.9,0.9}{%
	\makecell[{{p{\linewidth}}}]{% 
	  \tt {\tiny [A$\langle$GM]}  
	 ARGUMENT: {'Durch das Entfernen von A89 verlieren wir erheblichen Wert. Stattdessen könnten wir einen kleineren Gegenstand wie B85 (319) opfern, um Platz für B21 und A04 zu schaffen, während wir den hohen Wert von A89 behalten.'}\\ \tt \\ \tt VORSCHLAG: {'A41', 'A45', 'B61', 'A89', 'C64', 'A03', 'B21', 'A04'} 
	  } 
	   } 
	   } 
	 & \\ 
 

    \theutterance \stepcounter{utterance}  

    & \multicolumn{4}{p{0.6\linewidth}}{\cellcolor[rgb]{0.8,1,0.9}{%
	\makecell[{{p{\linewidth}}}]{% 
	  \tt {\tiny [A$\rangle$GM]}  
	 STRATEGISCHE ÜBERLEGUNGEN: {'Der Gegner akzeptiert nun B21 und A04, was gut ist, aber möchte B85 entfernen, das für mich einen guten Wert hat (461/319=1.45). Der neue Vorschlag: A41=140, A45=516, B61=448, A89=742, C64=77, A03=694, B21=366, A04=665. Summe=3648. Dies ist unter Budget (4026-3648=378 verbleibend). Ich könnte C49 (513) hinzufügen, aber das würde 4161 ergeben - noch unter Budget. C49 hat aber nur mäßigen Wert für mich (262). Besser wäre, B85 (319) hinzuzufügen, was 3967 ergäbe - noch unter Budget. Das wäre ein guter Kompromiss.'}\\ \tt \\ \tt ARGUMENT: {'Ihr aktueller Vorschlag ist gut, aber wir haben noch Budget übrig. B85 wäre eine ausgezeichnete Ergänzung - es bietet guten Wert bei moderaten Kosten und würde unser Budget optimal nutzen, ohne es zu überschreiten.'}\\ \tt \\ \tt VORSCHLAG: {'A41', 'A45', 'B61', 'A89', 'C64', 'A03', 'B21', 'A04', 'B85'} 
	  } 
	   } 
	   } 
	 & & \\ 
 

    \theutterance \stepcounter{utterance}  

    & & \multicolumn{4}{p{0.6\linewidth}}{\cellcolor[rgb]{0.9,0.9,0.9}{%
	\makecell[{{p{\linewidth}}}]{% 
	  \tt {\tiny [GM$\rangle$B]}  
	 ARGUMENT: {'Ihr aktueller Vorschlag ist gut, aber wir haben noch Budget übrig. B85 wäre eine ausgezeichnete Ergänzung - es bietet guten Wert bei moderaten Kosten und würde unser Budget optimal nutzen, ohne es zu überschreiten.'}\\ \tt \\ \tt VORSCHLAG: {'A41', 'A45', 'B61', 'A89', 'C64', 'A03', 'B21', 'A04', 'B85'} 
	  } 
	   } 
	   } 
	 & \\ 
 

    \theutterance \stepcounter{utterance}  

    & & & \multicolumn{4}{p{0.6\linewidth}}{\cellcolor[rgb]{1,0.85,0.72}{%
	\makecell[{{p{\linewidth}}}]{% 
	  \tt {\tiny [GM$\langle$B]}  
	 STRATEGISCHE ÜBERLEGUNGEN: {'Der Gegner schlägt vor, B85 (319 Aufwand, 668 Wert für mich) hinzuzufügen. Der aktuelle Vorschlag hatte einen Aufwand von 3738 (A41=140, A45=516, B61=448, A89=742, C64=77, A03=694, B21=366, A04=665, Summe=140+516+448+742+77+694+366+665=3648? Vorheriger Aufwand war 3738, aber ich muss neu rechnen: 140+516+448+742+77+694+366+665 = 3648, nicht 3738. Mit B85 (+319) = 3967, was unter 4026 liegt. Mein Wert würde von 3822 auf 3822 + 668 = 4490 steigen. Das ist ausgezeichnet für mich! Der Gegner scheint B85 nicht hoch zu bewerten, da er es einfach hinzufügt. Ich sollte sofort zustimmen, da dieser Deal meinen Wert maximiert und unter dem Budget bleibt.'}\\ \tt \\ \tt ARGUMENT: {'Das ist ein ausgezeichneter Vorschlag! Durch die Hinzunahme von B85 nutzen wir das Budget optimal und schaffen einen hohen Gesamtnutzen für beide Seiten.'}\\ \tt \\ \tt ZUSTIMMUNG: {'A41', 'A45', 'B61', 'A89', 'C64', 'A03', 'B21', 'A04', 'B85'} 
	  } 
	   } 
	   } 
	  \\ 
 

    \theutterance \stepcounter{utterance}  

    & & & \multicolumn{2}{p{0.3\linewidth}}{\cellcolor[rgb]{0.95,0.95,0.95}{%
	\makecell[{{p{\linewidth}}}]{% 
	  \tt {\tiny [GM$|$GM]}  
	 game successful 
	  } 
	   } 
	   } 
	 & & \\ 
 

    \theutterance \stepcounter{utterance}  

    & & & \multicolumn{2}{p{0.3\linewidth}}{\cellcolor[rgb]{0.95,0.95,0.95}{%
	\makecell[{{p{\linewidth}}}]{% 
	  \tt {\tiny [GM$|$GM]}  
	 end game 
	  } 
	   } 
	   } 
	 & & \\ 
 

\end{supertabular}
}

\end{document}
